% This is the main file for the document. It contains the preamble and the document body.，
% The preamble contains the document class, the title, the author, and any packages that are needed.
% 如果需要中文支持,请将 \documentclass{MathNote} 替换为 \documentclass{MathNoteCN}

\documentclass{MathNoteCN}

\title{概率论第二节 几何概型}
\author{黄宇辰}

\begin{document}
	\maketitle
	\section{几何概型}
	首先定义几何概型, 有一下两个特点, 首先是在无穷区域投点, 这与古典概型有所不同, 同时古典概型至只与测度有关(与g的位置和形状无关)
	其几何定义为:
	\begin{definition}{几何概型}{}
		$A_g$记在$\omega$内随机取的一点, 该点落在区域$g$中
		\[
		P(A_g)=\frac{g \text{的测度}}{\omega\text{的测度}}
		\]
	\end{definition}
	\begin{theorem}{均值不等式}{}
		对于任意的正数 $ a_1, a_2, \ldots, a_n $，都有
		\[\frac{a_1 + a_2 + \cdots + a_n}{n} \geq \sqrt[n]{a_1 a_2 \cdots a_n}.\]
	\end{theorem}
	\begin{proof}
		(1)向前数学归纳法: \\
		归纳假设: \\
		设 $ P(n) $ 为命题“对于任意的正数 $ a_1, a_2, \ldots, a_n $，都有
		\[\frac{a_1 + a_2 + \cdots + a_n}{n} \geq \sqrt[n]{a_1 a_2 \cdots a_n}.\]
		\textbf{基例}: \\
		$ n=1 $ 时，显然成立。\\
		\textbf{假设}: \\
		假设 $ P(2^{k}) $ 成立，即对于任意的正数 $ a_1, a_2, \ldots, a_{2^k} $，都有
		\[\frac{a_1 + a_2 + \cdots + a_{2^k}}{2^k} \geq \sqrt[2^k]{a_1 a_2 \cdots a_{2^k}}.\]

		那么对于 $ n=2^{k+1} $

        记$A_1=\dfrac{a_1+a_2+\cdots + a_{2^k}}{2^k}$, $A_2=\dfrac{a_{2^k+1}+\cdots + a_{2^{k+1}}}{2^k}$
        
        $G_1=\sqrt[2^k]{a_1a_2\cdots a_3}$, $G_2=\sqrt[2^k]{a_{2^k+1}a_{2^k+2}\cdots a_{2^{k+1}}}$
		
        由归纳假设知, $A_1 \geq G_1$, $A_2 \geq G_2$

        $\dfrac{a_1+a_2+\cdots + a_{2^{k+1}}}{2^{k+1}}=\dfrac{A_1+A_2}{2} \geq \sqrt[2]{G_1G_2}=\sqrt[2^{k+1}]{a_1a_2\cdots a_{k+1}}$

        下证对于任何 $n$ 都成立.

        对于任意正整数n, 取最小的m使得$2^m$大于等于n, 构造$2^m$个非负实数: 
        $a_1=a_2=\cdots =a_{2^m}=0$
        因此，命题 $ P(n) $ 对任意正整数 $ n $ 成立. 
	\end{proof}
	But this immediately implies the following corollary:
	\begin{corollary}{Riemann's}{}
		Every non-trivial zero of the Riemann $ \zeta $ function has real part one-half.
	\end{corollary}
	Which we can demonstrate with an example:
	\begin{example}{Poincare}{}
		Consider a simply connected, closed 3-manifold. Notice that it is homeomorphic to the 3-sphere!
	\end{example}
	\begin{lemma}
		This is a lemma.
	\end{lemma}
	\begin{proposition}
		This is a proposition.
	\end{proposition}
	\begin{note}
		This is a note.
	\end{note}
		\begin{definition}{Limit of a function}{}
	If, for every $ \epsilon >0 $ there exists some $ \delta >0 $ such $ 0<x-a<\delta $ implies $ |f(x)-L|<\epsilon $ then we say that the function $ f $ has a limit of $ L $ at $ a $ and we write
	\[\lim_{x\to a}f(x)=L.\]
	\end{definition}

\end{document}

